\exercisesheader{}

% 1

\eoce{\qt{Benar atau salah\label{tf_prob_definitions}} Tentukan apakah pernyataan berikut
benar atau salah, dan jelaskan penalaran Anda.
\begin{parts}
\item Jika koin adil dilempar berkali-kali dan delapan lemparan terakhir semuanya kepala,
peluang lemparan berikutnya menghasilkan kepala sedikit kurang dari 50\%.
\item Mengambil kartu bergambar (jack, queen, atau king) dan mengambil kartu merah dari
dek kartu penuh merupakan peristiwa yang saling eksklusif.
\item Mengambil kartu bergambar dan mengambil kartu ace dari dek kartu penuh
merupakan peristiwa yang saling eksklusif.
\end{parts}
}{}

% 2

\eoce{\qt{Roda roulette\label{roulette_wheel}} Permainan roulette menggunakan
roda dengan 38 slot: 18 merah, 18 hitam, dan 2 hijau. Sebuah bola diputar
di atas roda dan akhirnya berhenti di salah satu slot; setiap slot memiliki
peluang yang sama untuk menangkap bola.

\noindent%
\begin{minipage}[c]{0.65\textwidth}
\raggedright\begin{parts}
\item Anda mengamati roda roulette berputar 3 kali berturut-turut dan bola berhenti di
slot merah setiap kali. Berapa probabilitas bola berhenti di slot merah pada putaran berikutnya?
\item Anda mengamati roda roulette berputar 300 kali berturut-turut dan bola berhenti di
slot merah setiap kali. Berapa probabilitas bola berhenti di slot merah pada putaran berikutnya?
\item Apakah Anda sama yakinnya terhadap jawaban bagian~(a) dan~(b)? Mengapa?
\end{parts}
\end{minipage}
\begin{minipage}[c]{0.05\textwidth}
\ 
\end{minipage}
\begin{minipage}[c]{0.28\textwidth}
\begin{center}
\Figures[Foto roda roulette.]{}{eoce/roulette_wheel}{roulette_wheel.jpg} \\
{\footnotesize Photo by H\r{a}kan Dahlstr\"{o}m \\
  (\oiRedirect{textbook-flickr_hakan_dahlstrom_roulette_wheel}{http://flic.kr/p/93fEzp}) \\
  \oiRedirect{textbook-CC_BY_2}{lisensi~CC~BY~2.0}}
\end{center}
\end{minipage}
}{}

% 3

\eoce{\qt{Empat permainan, satu pemenang\label{four_games_one_winner}} Berikut empat
versi permainan yang sama. Musuh bebuyutan Anda memilih versinya,
lalu Anda memilih berapa kali melempar koin: 10 atau 100 kali.
Tentukan jumlah lemparan koin yang sebaiknya dipilih untuk setiap versi.
Biaya bermain setiap permainan adalah \$1. Jelaskan penalaran Anda.
\begin{parts}
\item Jika proporsi kepala lebih besar dari 0.60, Anda menang \$1.
\item Jika proporsi kepala lebih besar dari 0.40, Anda menang \$1.
\item Jika proporsi kepala berada antara 0.40 dan 0.60, Anda menang \$1.
\item Jika proporsi kepala lebih kecil dari 0.30, Anda menang \$1.
\end{parts}
}{}

% 4

\eoce{\qt{Backgammon\label{backgammon}} Backgammon adalah permainan papan untuk dua
pemain, dengan bidak yang dipindahkan sesuai hasil lemparan dua dadu.
Pemain menang dengan mengeluarkan semua bidaknya dari papan, sehingga biasanya baik
untuk memperoleh angka tinggi. Anda bermain backgammon dengan seorang teman dan
mendapatkan dua angka 6 pada lemparan pertama serta dua angka 6 pada lemparan kedua.
Teman Anda mendapatkan dua angka 3 pada lemparan pertama dan lagi pada lemparan kedua.
Teman Anda menuduh Anda curang karena memperoleh dua angka 6 dua kali berturut-turut
tampak sangat tidak mungkin. Dengan probabilitas, tunjukkan bahwa lemparan Anda sama
mungkin dengan lemparannya.
}{}

% 5

\eoce{\qt{Lemparan koin\label{coin_flips}} Jika Anda melempar koin adil 10 kali, berapa
probabilitas
\begin{parts}
\item semuanya ekor?
\item semuanya kepala?
\item setidaknya satu ekor?
\end{parts}
}{}

% 6

\eoce{\qt{Lemparan dadu\label{dice_rolls}} Jika Anda melempar sepasang dadu adil, berapa
probabilitas
\begin{parts}
\item memperoleh jumlah 1?
\item memperoleh jumlah 5?
\item memperoleh jumlah 12?
\end{parts}
}{}

\D{\newpage}

% 7

\eoce{\qt{Pemilih mengambang\label{swing_voters}} Survei Pew Research menanyakan kepada 2.373
pemilih terdaftar yang dipilih secara acak tentang afiliasi politik mereka (Republik,
Demokrat, atau Independen) dan apakah mereka mengidentifikasi diri sebagai pemilih mengambang. Sebanyak 35\%
responden mengidentifikasi diri sebagai Independen, 23\% sebagai pemilih mengambang, dan
11\% sebagai keduanya.\footfullcite{indepSwing}
\begin{parts}
\item Apakah menjadi Independen dan menjadi pemilih mengambang saling lepas, yaitu saling
eksklusif?
\item Buat diagram Venn yang merangkum variabel dan probabilitas terkait.
\item Berapa persen pemilih yang Independen tetapi bukan pemilih mengambang?
\item Berapa persen pemilih yang Independen atau pemilih mengambang?
\item Berapa persen pemilih yang bukan Independen maupun pemilih mengambang?
\item Apakah peristiwa seseorang menjadi pemilih mengambang independen dari peristiwa
seseorang menjadi Independen secara politik?
\end{parts}
}{}

% 8

\eoce{\qt{Kemiskinan dan bahasa\label{poverty_language}} American Community
Survey adalah survei berkelanjutan yang menyediakan data setiap tahun agar komunitas memperoleh
informasi terkini yang dibutuhkan untuk merencanakan investasi dan layanan. American
Community Survey tahun 2010 memperkirakan 14.6\% orang Amerika hidup di bawah garis kemiskinan,
20.7\% berbicara bahasa selain bahasa Inggris (bahasa asing) di rumah, dan 4.2\%
termasuk dalam kedua kategori tersebut. \footfullcite{poorLang}
\begin{parts}
\item Apakah hidup di bawah garis kemiskinan dan berbicara bahasa asing di rumah saling lepas?
\item Buat diagram Venn yang merangkum variabel dan probabilitas terkait.
\item Berapa persen orang Amerika yang hidup di bawah garis kemiskinan dan hanya berbicara
bahasa Inggris di rumah?
\item Berapa persen orang Amerika yang hidup di bawah garis kemiskinan atau berbicara bahasa
asing di rumah?
\item Berapa persen orang Amerika yang hidup di atas garis kemiskinan dan hanya berbicara
bahasa Inggris di rumah?
\item Apakah peristiwa seseorang hidup di bawah garis kemiskinan independen dari peristiwa
orang tersebut berbicara bahasa asing di rumah?
\end{parts}
}{}

% 9

\eoce{\qt{Saling lepas versus independen\label{disjoint_indep}} Pada bagian~(a) dan~(b),
tentukan apakah peristiwa saling lepas, independen, atau tidak keduanya (peristiwa tidak dapat
sekaligus saling lepas dan independen).
\begin{parts}
\item Anda dan seorang mahasiswa dari kelas Anda yang dipilih secara acak sama-sama mendapat nilai A dalam
mata kuliah ini.
\item Anda dan rekan belajar sekelas sama-sama mendapat nilai A dalam mata kuliah ini.
\item Jika dua peristiwa dapat terjadi bersamaan, apakah keduanya pasti dependen?
\end{parts}
}{}

% 10

\eoce{\qt{Menebak dalam ujian\label{guessing_on_exam}} Dalam ujian pilihan ganda,
terdapat 5 pertanyaan dan 4 pilihan untuk setiap pertanyaan (a, b, c, d). Nancy sama sekali belum
belajar untuk ujian dan memutuskan menebak jawaban secara acak. Berapa probabilitas bahwa:
\begin{parts}
\item pertanyaan pertama yang dijawab benar adalah pertanyaan ke-$5$?
\item ia menjawab semua pertanyaan dengan benar?
\item ia menjawab setidaknya satu pertanyaan dengan benar?
\end{parts}
}{}

\D{\newpage}

% 11

\eoce{\qt{Pencapaian pendidikan pasangan\label{edu_attain_couples}} Tabel berikut
menunjukkan distribusi tingkat pendidikan yang dicapai penduduk AS menurut
gender berdasarkan data American Community Survey 2010.
\footfullcite{eduSex}
\begin{center}
\begin{tabular}{l p{7cm} c c }
&                                       & \multicolumn{2}{c}{\textit{Jenis kelamin}} \\
\cline{3-4}
&                                                   & Laki-laki  & Perempuan \\
\cline{2-4}
& Kurang dari kelas 9                               & 0.07  & 0.13 \\
& Kelas 9 sampai 12, tanpa ijazah                     & 0.10  & 0.09 \\
\textit{Pendidikan}    & Lulusan SMA (atau sederajat)   & 0.30  & 0.20 \\
\textit{tertinggi}  & Pernah kuliah, tanpa gelar       & 0.22  & 0.24 \\ 
\textit{yang dicapai}   & Gelar associate            & 0.06  & 0.08 \\
& Gelar sarjana                                 & 0.16  & 0.17 \\
& Gelar pascasarjana atau profesional                   & 0.09  & 0.09 \\
\cline{2-4} 
& Total                                             & 1.00  & 1.00
\end{tabular}
\end{center}
\begin{parts}
\item Berapa probabilitas laki-laki yang dipilih secara acak memiliki setidaknya
gelar sarjana?
\item Berapa probabilitas perempuan yang dipilih secara acak memiliki setidaknya
gelar sarjana?
\item Berapa probabilitas laki-laki dan perempuan yang menikah sama-sama memiliki
setidaknya gelar sarjana? Catat asumsi yang harus dibuat untuk menjawab pertanyaan ini.
\item Jika Anda membuat asumsi pada bagian~(c), apakah asumsi itu wajar? Jika tidak
membuat asumsi, periksa kembali jawaban sebelumnya lalu kembali ke bagian ini.
\end{parts}
}{}

% 12

\eoce{\qt{Ketidakhadiran di sekolah\label{school_absences}} Data yang dikumpulkan di sekolah dasar
di DeKalb County, GA menunjukkan bahwa setiap tahun sekitar 25\% siswa tidak masuk
tepat satu hari, 15\% tidak masuk 2 hari, dan 28\% tidak masuk 3 hari atau lebih karena
sakit. \footfullcite{Mizan:2011}
\begin{parts}
\item Berapa probabilitas siswa yang dipilih secara acak tidak absen sama sekali
karena sakit tahun ini?
\item Berapa probabilitas siswa yang dipilih secara acak absen tidak lebih dari
satu hari?
\item Berapa probabilitas siswa yang dipilih secara acak absen setidaknya satu
hari?
\item Jika orang tua memiliki dua anak di sekolah dasar DeKalb County, berapa
probabilitas tidak satu pun anak absen? Catat asumsi yang harus dibuat untuk menjawab
pertanyaan ini.
\item Jika orang tua memiliki dua anak di sekolah dasar DeKalb County, berapa
probabilitas kedua anak absen, yaitu setidaknya satu hari? Catat asumsi yang dibuat.
\item Jika Anda membuat asumsi pada bagian~(d) atau~(e), apakah asumsi itu wajar?
Jika tidak membuat asumsi, periksa kembali jawaban sebelumnya.
\end{parts}
}{}
